\section{Strong local density limits and forward arithmetic likelihood}
\label{sec:v54-forward-likelihood}
Higher integrability upgrades the complete raw local law and the
positive physical remainders. It also distinguishes two questions
which were both left open by the lower local law: forward relative
entropy can now be controlled, although forward essential likelihood
and the two-sided pointwise local theorem still require physical
height estimates.

Retain $q_*=1+s_*$ from \eqref{eq:v54-admissible-exponent}, and fix
$1<q<q_*$. All local integrals below concern a translated interval
$I_t=[t,t+h]$ of fixed length $h>0$. The normalized original density
and signed transition kernel are $P$ and $G$ as in the preceding
section. For this proof, the elementary interpolation inequality is
\begin{equation}\label{eq:v54-local-interpolation}
 \int_I |v|^q
 \le\left(\int_I |v|\right)^{\theta_q}
       \left(\int_I |v|^{q_*}\right)^{1-\theta_q},
 \qquad \theta_q=\frac{q_*-q}{q_*-1}\in(0,1).
\end{equation}
It follows from H\"older's inequality applied to
$|v|^q=|v|^{\theta_q}(|v|^{q_*})^{1-\theta_q}$.

\subsection{Physical remainders in a stronger local norm}
\begin{theorem}[Higher-integrable control of both physical sources]
\label{thm:v54-physical-Lq}
For the positive physical remainder of
\eqref{eq:v48-physical-remainder}, put
$B_\varepsilon=m^2b_{n,k,m,R}^{\varepsilon,1}$. Uniformly for
$m\ge2$ and $0<\varepsilon<\varepsilon_0$,
\begin{equation}\label{eq:v54-physical-Lq}
 \sup_{R,n,k,t}\int_{I_t}B_\varepsilon^q
 \le C_q(1+h)
              (\varepsilon^{1/16}+e^{-\kappa m})^{\theta_q}.
\end{equation}
The same bound holds for each first-incidence and first-clearance
component separately. For an arbitrary measurable insertion $|w|\le M$
it holds with $|m^2b^{\varepsilon,w}|^q$ and a factor $M^q$.
Consequently
\begin{align}
 &\limsup_{m\to\infty}\sup_{R,n,k,t,\,|w|\le M,\,B>0}
 \int_{I_t}\left|m^2\bigl(b^{\varepsilon,w}
                       -K_B*b^{\varepsilon,w}\bigr)(u)\right|^qdu
 \le C_q M^q(1+h)\varepsilon^{\theta_q/16}.
 \label{eq:v54-physical-all-band-Lq}
\end{align}
This is not an essential-height estimate.
\end{theorem}
\begin{proof}
The physical measure is a positive submeasure of the original source,
so $0\le B_\varepsilon\le P$ almost everywhere. Its local mass is
bounded by \eqref{eq:v49-physical-window-bound}, at the same fixed
auxiliary band as above, by
$C(1+h)(\varepsilon^{1/16}+e^{-\kappa m})$.
Theorem~\ref{thm:v54-uniform-integrability} bounds its $q_*$ moment.
Apply \eqref{eq:v54-local-interpolation}. Each disjoint physical type
is dominated by their sum, and Radon--Nikodym domination gives
$|b^{\varepsilon,w}|\le M b^{\varepsilon,1}$.

For completeness, Minkowski's integral inequality gives
\[
 \sup_t\|K_B*v\|_{L^q(I_t)}
       \le\|K_1\|_1\sup_t\|v\|_{L^q(I_t)}.
\]
Apply this finite inequality before taking any supremum or limit.
It holds also for the real sign-changing reconstruction kernel. The
factor $1+\|K_1\|_1$ is independent of $B$, which proves
\eqref{eq:v54-physical-all-band-Lq}. No bandwidth depending on $m$
has entered a spectral estimate.
\end{proof}

\begin{theorem}[The complete arithmetic raw law in local $L^q$]
\label{thm:v54-raw-Lq}
For every fixed $h>0$ and $1<q<q_*$,
\begin{equation}\label{eq:v54-raw-Lq}
 \sup_{R,\,1\le n\le m,\,k,t}\frac1h\int_{I_t}
  \left|m^2p_{n,R}(k,m,u)-\mathcal L_{m,R}(k_1,k_2,u,n)\right|^qdu
                         \longrightarrow0.
\end{equation}
For the pinned collision-path measure $\mathbf P(u)$ and Gaussian
bridge law $\mathsf W_R$, the same limit holds with the integrand
\begin{equation}\label{eq:v54-path-Lq}
 \|\mathbf P(u)-G(u)\mathsf W_R\|_{\mathrm{BL}^*}^{q}.
\end{equation}
For the actual-return path this conclusion holds on the central
classes of Lemma~\ref{lem:v52-return-transfer}. The path metric and
bounded-Lipschitz convention are unchanged.
\end{theorem}
\begin{proof}
The scalar $L^1$ error tends to zero uniformly by
Theorem~\ref{thm:v49-raw-local-TV}. The kernel $G$ is uniformly
bounded, as in the finite moving-peak formula. Hence
\[
 \int_{I_t}|P-G|^{q_*}\le C_*(1+h).
\]
Interpolation proves \eqref{eq:v54-raw-Lq}; division by the fixed
$h$ is harmless. This gives a modulus in the already proved local
$L^1$ error, not a rate in $m$.

For the path statement put
$v(u)=\|\mathbf P(u)-G(u)\mathsf W_R\|_{\mathrm{BL}^*}$.
The constant function one belongs to the test ball, a positive
path measure has this norm equal to its mass, and thus
$v(u)\le P(u)+|G(u)|$. Its $q_*$ moment is bounded by the same
scalar estimate, while its $L^1$ norm tends to zero by
Theorem~\ref{thm:v52-path-remainder}. Apply the same interpolation.
The countable norming class makes $v$ measurable off the common
roof-null set. For the return path use the roof-integrated transfer
in Lemma~\ref{lem:v52-return-transfer}; its positive path measure
again has mass $P$. No path-space total variation is asserted.
\end{proof}

\subsection{Forward likelihood and relative entropy}
Fix $d>0$ and restrict to windows $I_t$ on which $G\ge d$ almost
everywhere. This is a pointwise reference-floor condition, not merely
a lower bound for its integrated mass. Put
\[
 F=\int_{I_t}P(u)\,du,\qquad H=\int_{I_t}G(u)\,du,
 \qquad d\mathbb P=\one_{I_t}\frac{P}{F}\,du,
 \qquad d\mathbb Q=\one_{I_t}\frac{G}{H}\,du.
\]
Here $\mathbb P$ is precisely the roof law under the unchanged event
$K_{n,R}=k$, $N_{n,R}=m$, $T_{n,R}\in I_t$.
The factors $m^2$ cancel in its normalization. The reference $G$ is
positive only because of the stated floor on this window.

\begin{theorem}[Forward finite-order arithmetic likelihood]
\label{thm:v54-forward-likelihood}
Uniformly over these windows and exact targets, for every $1<q<q_*$,
\begin{equation}\label{eq:v54-likelihood-Lq}
 \int\left|\frac{d\mathbb P}{d\mathbb Q}-1\right|^q
                                                   d\mathbb Q\to0.
\end{equation}
The forward relative entropy and every fixed forward R\'enyi
divergence of order $1<s<q_*$ consequently tend to zero:
\begin{equation}\label{eq:v54-forward-divergences}
 D(\mathbb P\Vert\mathbb Q)\to0,\qquad
 D_s(\mathbb P\Vert\mathbb Q)
 =\frac1{s-1}\log\int
        \left(\frac{d\mathbb P}{d\mathbb Q}\right)^s d\mathbb Q\to0.
\end{equation}
Combined with Theorem~\ref{thm:v51-conditional-minorization}, this
also proves convergence of the sum of the two directed relative
entropies. A forward essential likelihood bound is not asserted.
\end{theorem}
\begin{proof}
The local law gives $F-H\to0$ uniformly, $H\ge dh$, and $F\ge dh/2$
for sufficiently large counts. Since $G$ is uniformly bounded,
$F,H$ have uniform upper bounds on the fixed window. Thus
$\mathbb P\ll\mathbb Q$, with likelihood
$r(u)=HP(u)/(FG(u))$. For a constant depending only on $d,h,q$,
\[
 \int|r-1|^q\,d\mathbb Q
 =\frac1H\int_{I_t}G^{1-q}|(H/F)P-G|^q\,du
 \le C_{d,h,q}\left(\int_{I_t}|P-G|^q\,du+|F-H|^q\right).
\]
Theorem~\ref{thm:v54-raw-Lq} proves \eqref{eq:v54-likelihood-Lq}.
For $1<s<q_*$ apply it with $q=s$. The $L^s(\mathbb Q)$ triangle
inequality gives $|\|r\|_s-1|\le\|r-1\|_s$, proving the R\'enyi
assertion.

To verify the entropy step without assuming uniform boundedness of
$r$, choose any $1<q<q_*$; in particular $q<2$. For $x\ge0$,
\[
 0\le x\log x-x+1\le C_q|x-1|^q,
 \qquad 0\log0=0.
\]
Taylor's formula proves this on $[1/2,2]$ since the quadratic error
is bounded by a constant times $|x-1|^q$ there. On $[0,1/2]$ use
boundedness and $|x-1|\ge1/2$; on $[2,\infty)$ use
$x\log x\le C_q x^q$. Integrate with $x=r$ and
$\int(r-1)d\mathbb Q=0$. This proves
$D(\mathbb P\Vert\mathbb Q)\le C_q\|r-1\|_q^q\to0$.
The reverse entropy already tends to zero by conditional minorization.
No estimate of $\operatorname*{ess\,sup}r$ is used.
\end{proof}

\begin{remark}[Arithmetic and positive reference classes]
At fixed $R$, on a central compact set and a positive arithmetic
class $\mathfrak a_R(k,n,m)\ge a_0>0$, the fixed-radius kernel
specialization supplies the reference floor for sufficiently large
$m$, as in Remark~\ref{rem:v51-arithmetic-positivity}. Uniformly
through changes of arithmetic type the floor on $\mathcal L_{m,R}$
is retained explicitly. A zero residue is not assigned a positive
conditional reference. The roof interval and the discrete labels
are identical on both sides of every likelihood comparison.
\end{remark}

\paragraph{The remaining height question.}
Higher integrability is a stronger scalar conclusion than local total
variation, and forward entropy does not follow from the old reverse
minorization alone. Nevertheless finite $L^q$ control does not imply
an essential-supremum bound. For example, with the present fixed $q_*>1$,
normalize $1+j\one_{[0,j^{-(q_*+1)}]}$ to a probability density on
$[0,1]$. Its $L^{q_*}$ norms stay bounded, its forward entropy relative
to the constant density tends to zero, and its heights tend to
infinity. This only distinguishes the topologies; it is not a billiard
counterexample. The positive incidence and clearance heights in
\eqref{eq:v51-positive-criterion}, and hence the two-sided pointwise
raw law and unconditional uniform same-roof bridge, remain to be
estimated on the actual source.
